% CS466 Deep Learning - Term Project Proposal
% Compile with pdfLaTeX (Overleaf default).

\documentclass[11pt]{article}

\usepackage[utf8]{inputenc}
\usepackage[T1]{fontenc}
\usepackage{lmodern}
\usepackage[margin=0.9in]{geometry}
\usepackage{amsmath}
\usepackage{booktabs}
\usepackage{enumitem}
\usepackage[hidelinks]{hyperref}

\setlist{itemsep=1pt, topsep=3pt}

\title{\textbf{Small and Fast Depth Estimation with Few Labels\\by Learning from a Big Model}\\[4pt]
\large CS466 Deep Learning, Fall 2026 --- Term Project Proposal (Models Track)}
\author{Oğuz Kağan Göçmez\\
\small Özyeğin University\\
\small \texttt{kagan.gocmez@ozu.edu.tr}}
\date{October 30, 2026}

\begin{document}
\maketitle

%--------------------------------------------------------------------
\section{Introduction}

A robot needs to know how far away things are. Without this information,
it cannot avoid obstacles or plan where to go. Depth sensors such as LiDAR
can measure distance, but they are expensive. A normal camera is cheap and is
already on almost every robot, so it would be very useful if a model could
estimate depth from just one camera image. This task is called
\emph{monocular depth estimation}~\cite{eigen2014}.

Today there are very strong models for this task, such as Depth Anything
V2~\cite{yang2024dav2}. However, they are big and slow, so they do not run
well on small robots or phones. Small models are fast, but they are less
accurate. There is also a second problem: training a depth model needs real
depth maps, and collecting them needs a depth sensor and a lot of effort.

In this project, I will train a small model (the \emph{student}) that learns
from a big model (the \emph{teacher}). This idea is called \emph{knowledge
distillation}~\cite{hinton2015}. The teacher can label any image for free, so
the student can also learn from images that have no real depth.

Distillation for depth has been studied before~\cite{wang2021kd}, and
ZipDepth~\cite{tosi2026zipdepth} recently distilled a foundation model using
14.1M images and several GPUs. My project looks at the opposite end:
\textbf{my contribution is a distillation recipe for very small students
(1--4M parameters) when only a small part of the training images has real
depth, together with a careful study of when the teacher helps and when it
does not.} My first feasibility test (Section~\ref{sec:prelim}) showed that
the teacher does not help when every image has real depth, which is why I
focus on the few-label setting.

%--------------------------------------------------------------------
\section{Problem Statement}

\paragraph{Task.} Given one RGB image, predict a depth map, i.e.\ a distance
value in meters for every pixel.

\paragraph{Dataset.} I will use NYU Depth V2~\cite{silberman2012}, which
contains indoor scenes recorded with a Microsoft Kinect. I use the common
training subset of 50,688 images from 284 scenes~\cite{alhashim2018} and the
standard test set of 654 images. To simulate few labels, I keep the real depth
for only a part of the training images (1\%, 5\%, 25\% or 100\%) and treat the
rest as unlabeled. Frames from the same room look very similar, so I will hold
out a validation set \emph{by scene} and use it for all tuning; the 654 test
images are used only for the final results.

\paragraph{Expected results.} I expect that:
\begin{itemize}
  \item with all labels, the teacher adds little (my first test already shows this);
  \item with few labels (1--5\%), a student that also learns from the teacher's
        outputs on the unlabeled images will be clearly more accurate than a
        student trained only on the labeled images;
  \item this gap will shrink as the labeled part grows;
  \item the student will be about 8 times smaller than the teacher and run in
        real time on a laptop GPU.
\end{itemize}

\paragraph{Evaluation.} I will use the standard depth metrics: absolute
relative error (AbsRel), root mean squared error (RMSE) and threshold accuracy
$\delta < 1.25$. To measure efficiency, I will report the number of
parameters, FLOPs and inference time. I will also compare with published
lightweight models such as FastDepth~\cite{wofk2019}.

%--------------------------------------------------------------------
\section{Technical Approach}

\paragraph{Teacher.} Depth Anything V2 Small~\cite{yang2024dav2} (24.8M
parameters). I run it only once on all training images and save its outputs,
so the teacher is not needed during training. Its output is relative inverse
depth, which is correct only up to an unknown scale and shift. On labeled
images this could be fixed with the real depth, but unlabeled images have no
real depth. So for the teacher term I use a \emph{scale-and-shift-invariant}
loss~\cite{ranftl2022}: for each image, the student's prediction (converted to
inverse depth) is first aligned to the teacher output by least squares, and
only then compared. This way, the teacher term never uses real depth.

\paragraph{Student.} A small encoder--decoder network: a pretrained
MobileNetV3~\cite{howard2019} or EfficientNet-B0~\cite{tan2019} encoder and a
light decoder similar to FastDepth~\cite{wofk2019}. I use three sizes (about
1M, 3M and 4M parameters) and an input size of $256 \times 192$.

\paragraph{Training loss.}
\begin{equation}
  \mathcal{L} \;=\; \lambda_{1}\,\mathcal{L}_{\text{gt}}
  \;+\; \lambda_{2}\,\mathcal{L}_{\text{teacher}}
  \;+\; \lambda_{3}\,\mathcal{L}_{\text{edge}}
\end{equation}
$\mathcal{L}_{\text{gt}}$ is the scale-invariant log loss~\cite{eigen2014}
against the real depth, used only on labeled images; it teaches the student
the real scale in meters. $\mathcal{L}_{\text{teacher}}$ is the
scale-and-shift-invariant L1 loss against the teacher, used on all images; it
teaches the shape of the scene. $\mathcal{L}_{\text{edge}}$ matches the image
gradients of the aligned prediction and the teacher output, so the student
learns sharp object borders.

\paragraph{Experiments.} I train the same student several times and change
only one thing each time. The main setting uses 5\% labels.

\begin{center}
\small
\begin{tabular}{@{}cl@{}}
\toprule
\textbf{ID} & \textbf{Setting} \\
\midrule
E1 & Real depth only, on the labeled images (baseline) \\
E2 & Real depth + teacher, on the labeled images only (does the teacher signal help by itself?) \\
E3 & Real depth on the labeled images + teacher on all images \\
E4 & E3 + edge loss (full recipe) \\
E5 & E1 vs.\ E4 with 1\% / 5\% / 25\% / 100\% labels \\
E6 & E4 with student sizes of about 1M / 3M / 4M parameters \\
\bottomrule
\end{tabular}
\end{center}

All runs use the same number of training steps (about 63k steps at batch size
8, roughly 10 epochs over 50K images), and I pick the best checkpoint on the
validation set, so the runs are fair to compare. I will also show depth maps
for success and failure cases.

%--------------------------------------------------------------------
\section{Preliminary Results}
\label{sec:prelim}

To check that the plan is realistic, I wrote a feasibility script and ran it on
my laptop (Apple M4 Pro, 24\,GB). It runs the teacher, builds the 3.15M student
and trains it for 1,500 steps on 2,000 random training images (about 7
minutes), then tests on 200 test images:

\begin{center}
\small
\begin{tabular}{@{}lccc@{}}
\toprule
\textbf{Model} & \textbf{AbsRel}~$\downarrow$ & \textbf{RMSE}~$\downarrow$ & $\boldsymbol{\delta<1.25}$~$\uparrow$ \\
\midrule
Constant depth (sanity check) & 0.416 & 1.209 & 0.386 \\
Student, real depth only      & 0.202 & 0.675 & 0.675 \\
Student, real depth + teacher & 0.204 & 0.677 & 0.673 \\
Teacher, aligned with real depth (optimistic reference) & 0.093 & 0.430 & 0.911 \\
\bottomrule
\end{tabular}
\end{center}

The student clearly learns, but the teacher gives no gain here. I checked the
data: the NYU training depth maps are already filled, with 0\% missing pixels,
so with full labels the teacher mostly repeats information that the student
already has. This result led me to the few-label setting. An earlier test on a
copy of NYU that covers only a few rooms overfit badly ($\delta<1.25 = 0.33$),
which is why I will split the validation set by scene.

The teacher needs about 40\,ms per image (0.6 hours for all 50K images). The
student needs 0.74 GFLOPs and 8.6\,ms per image on the laptop GPU (peak memory
2.8\,GB). One full run takes about 4.6 hours, so the 12 runs of E1--E6 need
about 55 GPU hours, which I will split between my laptop and Colab Pro.

\paragraph{Plan.} \textbf{Nov 2--8:} teacher outputs, scene-based validation
split, teacher loss. \textbf{Nov 9--22:} E1 at all label budgets, E3 and E4 at
5\%. \textbf{Nov 23--Dec 13:} E2, E5, E6 and analysis. \textbf{Dec 14--27:}
figures, final report and presentation. If time becomes short, I will keep
only the 5\% and 100\% budgets in E5 and drop E6.

%--------------------------------------------------------------------
\begin{thebibliography}{99}
\footnotesize
\setlength{\itemsep}{0pt}

\bibitem{eigen2014}
D.~Eigen, C.~Puhrsch, and R.~Fergus.
Depth map prediction from a single image using a multi-scale deep network.
In \emph{NeurIPS}, 2014.

\bibitem{yang2024dav2}
L.~Yang, B.~Kang, Z.~Huang, Z.~Zhao, X.~Xu, J.~Feng, and H.~Zhao.
Depth Anything V2.
In \emph{NeurIPS}, 2024.

\bibitem{hinton2015}
G.~Hinton, O.~Vinyals, and J.~Dean.
Distilling the knowledge in a neural network.
\emph{arXiv:1503.02531}, 2015.

\bibitem{wang2021kd}
Y.~Wang, X.~Li, M.~Shi, K.~Xian, and Z.~Cao.
Knowledge distillation for fast and accurate monocular depth estimation on
mobile devices.
In \emph{CVPR Workshops}, 2021.

\bibitem{tosi2026zipdepth}
F.~Tosi, L.~Bartolomei, M.~Poggi, and S.~Mattoccia.
ZipDepth: Bringing lightweight zero-shot monocular depth anywhere, on any
device.
\emph{arXiv:2607.08771}, 2026.

\bibitem{silberman2012}
N.~Silberman, D.~Hoiem, P.~Kohli, and R.~Fergus.
Indoor segmentation and support inference from RGBD images.
In \emph{ECCV}, 2012.

\bibitem{alhashim2018}
I.~Alhashim and P.~Wonka.
High quality monocular depth estimation via transfer learning.
\emph{arXiv:1812.11941}, 2018.

\bibitem{ranftl2022}
R.~Ranftl, K.~Lasinger, D.~Hafner, K.~Schindler, and V.~Koltun.
Towards robust monocular depth estimation: Mixing datasets for zero-shot
cross-dataset transfer.
\emph{IEEE TPAMI}, 2022.

\bibitem{howard2019}
A.~Howard et~al.
Searching for MobileNetV3.
In \emph{ICCV}, 2019.

\bibitem{tan2019}
M.~Tan and Q.~V.~Le.
EfficientNet: Rethinking model scaling for convolutional neural networks.
In \emph{ICML}, 2019.

\bibitem{wofk2019}
D.~Wofk, F.~Ma, T.-J.~Yang, S.~Karaman, and V.~Sze.
FastDepth: Fast monocular depth estimation on embedded systems.
In \emph{ICRA}, 2019.

\end{thebibliography}

\end{document}
